\scanpage{370}\mspage{286}
\unsure{Ce|le} foncteur d'inclusion est accessible,
\unsure{en fait} il est $\pi$-accessible, puisque
(4.13.2.1) l'est (cf 4.12.3).

\uline{Définition} 4.13.3. Soient $M$, $S\subset\operatorname{Ob}M$, \inserted{$\underline S\subset M$,} comme
dans 4.12.1, \inserted{$\pi$ un cardinal infini.} On dit que $S$ est
une partie \uline{$\pi$-accessible} de $\operatorname{Ob}M$,
ou que $\underline S$ est une sous-catégorie
$\pi$-accessible dans $M$, si \deleted{$\pi\geq\pi_0$ \unsure{déf}}
\deleted{\unsure{1)|a)} $\pi$ est \unsure{quel|plus} cardinal
plus petit cardinal utile pour $M$,
et si de plus} \deleted{\unsure{b)}} $\unsure{S|C},\pi$ \unsure{satisfont|satisfait} les
conditions de 4.13.3 a), i.e.
$\underline S$ est l'image essentielle de $\mathrm{Ind}_\pi(S(\pi))$,
où $S(\pi)$ est une sous-catégorie pleine
de $M_\pi$.

Notons que, quitte à remplacer
$S(\pi)$ par $S'(\pi)=$ \deleted{sous-catégorie \unsure{st}}\inserted{\unsure{$\leftrightarrows$-plétion|complétion} de} $S(\pi)$
dans $M_\pi$, OPS que $S(\pi)$ est une
sous-catégorie \uline{str. pleine} \uline{et karoubienne}
% archive sheet number: 369
